% Pursuing Stacks (cotes 134-2 to 134-8): archive cote 134-7, site page 73 (PDF page 74), typescript page 515.
% typescript. Draft compared: ps6.tex:3648-3720.
% AI-transcribed by Claude Opus 5.5 (claude-opus-5-5, adaptive thinking, effort high), Message Batches API, batch ps4 (msgbatch_012Aoc1NuskATPsyBeKSYK7J), 2026-11-12; not yet checked by eye.
\scanpage{134-7/73}
\mspage{515}

topy theory.  \xout{Of} I was led to postulate the existence of a functor $\overset{\hand{$\sim$}}{S}$  in
the context of \uline{shhematic} homotopy types,  by reasons of symmetry in the
LH.  and L\hand{$\pi$}. formalism,  which I'll try to get through a little later.
It should be clear from the \unsure{outset|putset} though that in the schematic set-up,
even the mere existence of (schematic) homotopy types $S(\underline{Z},n)$  (correspon-%
ding to ordinary spheres),  and \xout{xx} even for $n=2$ or $3$\inserted{only}, is very far from
trivial,  and as a matter of fact is'nt even known (in the set-up of
semisimplicial unipotent bundles).  This shows at the same time that if
such a functor $\overset{\hand{$\sim$}}{S}$ can actually be constructed in the schematic case,  it is
likely to be a lot more interesting still than in the discrete case,  as
it will presumably give information on homotopy groups of spheres (via the
desciption of the \hand{$\pi$}$_i$'s  of a semisimplicial unipotent bundle via the Lie
functor  (cf. section 118)\hand{)}.

\margin{\unsure{+}}There  is a fourth tentative functor in between the categories
Hotab  and  Hot ,  or more accurately\unsure{;}  from a suitable subcategory
H\over{o}{\unsure{$\breve{o}$}}t$^\bullet$(0)  of Hot$^\bullet$  (corresponding to 0-connected pointed homotopy types
whose \hand{$\pi$}$_1$ is abelian and acts trivially on the higher \hand{$\pi$}$_i$'s) to
Hotab (0-conn) = H\over{o}{\unsure{$\hat{o}$}}tab(0),  strongly suggested by the schematic case of
section 118,  namely a functor

(*) \qquad $\mathrm{L}\hand{$\pi$}. : \mathrm{Hot}^\bullet(0) \hand{$\longrightarrow$} \mathrm{Hotab}^\bullet(0) \quad ???$  ,

whose main functorial property should be

(\over{**}{12}) \qquad $H_i(\mathrm{L}\hand{$\pi$}.(X)) \hand{$\simeq$} \hand{$\pi$}_i(X)$ \quad \deleted{(?)} .

To play really safe,  one might hope such a functor to be defined at any
rate for 1-connected pointed homotopy types.  Another equally important
property, suggested by the schematic set-up,\inserted[typed]{as \unsure{wehh} \unsure{ad} by (11)}\inserted{and (**),} is the formula

(\over{***}{13}) \qquad $\mathrm{L}\hand{$\pi$}.(\overset{\hand{$\sim$}}{K}(L.)) \hand{$\simeq$} L.$ \quad \deleted{(?)} ,

similar to formula (10)  (with the pair  ($\overset{\hand{$\sim$}}{S}$, $\overset{\hand{$\sim$}}{\mathrm{LH}}$. )  replaced by the pair
($\overset{\hand{$\sim$}}{K}$ , L\hand{$\pi$}. )).  Of course,  for X of the type $\overset{\hand{$\sim$}}{K}$(L.),  (\over{**}{12}) follows from
(\over{***}{13}) .

3.10.

\xout{Applying} We'll see though that a functor  L\hand{$\pi$}.  \xout{\unsure{(*)}} as in (*),  satis-%
fying the properties above,  \uline{does not exist}.  Indeed,  applying L\hand{$\pi$}. to
the adjunction morphism
\[ X \hand{$\longrightarrow$} \overset{\hand{$\sim$}}{K}\ \overset{\hand{$\sim$}}{\mathrm{LH}}.(X) \]
(cf. p.509, (6) ),\inserted{and using (13),} we should get a Hurewicz  homomorphism

\margin{73}
